\documentclass[a4paper]{article}
\usepackage[pdftex]{graphicx}
\usepackage[utf8]{inputenc}
\usepackage{enumerate}
\usepackage{amssymb}
\usepackage{tikz}
\usepackage{href-ul}
\hypersetup{
	colorlinks=true,
	linkcolor=blue,
	urlcolor=blue}
\usepackage{geometry}
\geometry{a4paper, top=15mm, left=15mm, right=15mm, bottom=15mm,
	headsep=10mm, footskip=12mm}

% Start the document
\begin{document}
\begin{tikzpicture}(10,3)
	\draw[ultra thick](2,0) --(10,0) -- (10,1.3) --(2,1.3) -- (2,0);
	\draw[fill=black](2,0)-- (0,.6) -- (2,.6) -- (2,0);
	\node at (6,.6) {\large Lösungsblatt von \href{https://www.okuyakl.de}{www.okuyakl.de}};
\end{tikzpicture}
\vspace{0.5 cm}

\noindent{\bf L Ö S U N G}\\
\renewcommand{\arraystretch}{3}
\begin{tabular}{|p{40 pt}|p{100 pt}|p{100 pt}|p{100 pt}|}
\hline
Funktion & Y-Achsenabschnitt $b$ &Wachstumsfaktor $a$  & Funktionsterm $y=$ \\
\hline
$f_1(x):$ &  4    & $\sqrt{2}$      &  $y=4\cdot(\sqrt{2})^x$ \\
\hline
$f_2(x):$ & 1     & $2$                &  $y=2^x$\\
\hline
$f_3(x):$ & 0,5 &  ${1\over 3}$ &  $y=0,5\cdot \left({1\over 3}\right)^x $\\
\hline
$f_4(x):$ & -1   & $4$                &  $y=-4^x$\\
\hline
$f_5(x):$ & -2   &  $\sqrt{1\over 2}$ &  $y=-2\cdot \left(\sqrt{1\over 2}\right)^x$\\
\hline
$f_6(x):$ & -4   &  $\sqrt[3]{2}$  & $y=-4 \cdot (\sqrt[3]{2})^x $ \\
\hline
\end{tabular}

\begin{center}
	\includegraphics[width=7 cm]{../../viecher/endcomic.pdf}
	
	Hier geht es zurück zum \href{https://www.okuyakl.de/math/m10efuL018/aa018.pdf}{Aufgabenblatt}
\end{center}


%Ende
\end{document}